\begin{table}[htbp]
\centering
\footnotesize
\caption{Problema D, meia-lua, Passos 2 a 5 (openai/gpt-4o-mini): métrica em R$^2$, R$^3$ e R$^4$ por regime.}
\label{tab:D_passos_2_a_5_gpt4mini}
\begin{tabular}{lrrrrlrrlr}
\toprule
 &  &  & \multicolumn{3}{c}{Perceptron} & \multicolumn{3}{c}{NNLS} &  \\ \cmidrule(lr){4-6} \cmidrule(lr){7-9}
Regime & Acur.\ LLM vs real & Partição LLM (n0/n1) & Fid.\ (vs LLM) & Acur.\ vs real & $\hat{W}$ unit. & Fid.\ (vs LLM) & Acur.\ vs real & $\hat{W}$ unit. & $\Delta$ fid.\ vs R$^2$ \\
\midrule
R$^2$ — zero-shot & — & — & — & — & — & — & — & — & — \\
R$^2$ — 4-shot† & 81,9 (68,5–92,5) & 83/63 & 89,8 (76,7–97,3) & 82,9 (76,7–86,7) & [0,36; 0,89] & 88,7 (77,4–97,3) & 83,6 (78,7–86,7) & [0,21; 0,97] & +0,0 \\
R$^2$ — 10-shot† & 89,2 (83,6–96,4) & 73/67 & 92,3 (81,4–97,1) & 85,1 (77,3–88,0) & [0,31; 0,94] & 91,5 (80,7–96,4) & 85,2 (82,7–86,7) & [0,26; 0,96] & +0,0 \\
R$^2$ — 20-shot† & 90,4 (86,9–95,4) & 68/62 & 93,8 (92,3–96,2) & 86,1 (83,3–88,0) & [0,26; 0,97] & 92,9 (90,0–96,2) & 85,6 (84,7–86,7) & [0,26; 0,96] & +0,0 \\
R$^2$ — 40-shot† & 91,6 (86,4–94,5) & 58/52 & 92,1 (88,2–94,5) & 86,1 (84,0–87,3) & [0,26; 0,97] & 92,0 (88,2–93,6) & 85,9 (84,7–87,3) & [0,24; 0,97] & +0,0 \\
R$^2$ — W ajustado ao rótulo real (teto)† & — & — & — & 86,9 (86,7–87,3) & [0,22; 0,97] & — & 86,7 & [0,22; 0,97] & — \\
\midrule
R$^3$ — zero-shot & — & — & — & — & — & — & — & — & — \\
R$^3$ — 4-shot† & 81,9 (68,5–92,5) & 83/63 & 90,3 (78,8–97,3) & 82,2 (78,0–86,7) & [0,34; 0,66; 0,64] & 90,6 (77,4–99,3) & 83,1 (78,7–86,0) & [0,13; 0,60; 0,52] & +0,5 \\
R$^3$ — 10-shot† & 89,2 (83,6–96,4) & 73/67 & 92,5 (81,4–97,1) & 85,3 (78,0–87,3) & [0,23; 0,69; 0,66] & 92,0 (80,7–96,4) & 84,9 (81,3–86,7) & [0,24; 0,86; 0,19] & +0,2 \\
R$^3$ — 20-shot† & 90,4 (86,9–95,4) & 68/62 & 93,5 (91,5–96,2) & 85,6 (82,7–87,3) & [0,22; 0,72; 0,65] & 92,6 (90,0–96,2) & 85,1 (81,3–86,7) & [0,24; 0,83; 0,32] & $-$0,3 \\
R$^3$ — 40-shot† & 91,6 (86,4–94,5) & 58/52 & 92,1 (90,0–94,5) & 85,6 (84,0–86,7) & [0,19; 0,72; 0,66] & 91,4 (87,3–93,6) & 85,4 (84,0–87,3) & [0,22; 0,84; 0,33] & +0,0 \\
R$^3$ — W ajustado ao rótulo real (teto)† & — & — & — & 87,1 (86,7–87,3) & [0,16; 0,69; 0,71] & — & 86,7 & [0,22; 0,97; 0,00] & — \\
\midrule
R$^4$ — zero-shot & — & — & — & — & — & — & — & — & — \\
R$^4$ — 4-shot† & 81,9 (68,5–92,5) & 83/63 & 94,2 (89,7–99,3) & 81,7 (75,3–86,7) & [0,16; 0,72; 0,14; 0,63] & 93,3 (87,7–99,3) & 83,3 (79,3–87,3) & [0,25; 0,87; 0,09; 0,12] & +4,4 \\
R$^4$ — 10-shot† & 89,2 (83,6–96,4) & 73/67 & 93,6 (80,0–98,6) & 85,7 (80,7–88,0) & [0,17; 0,74; 0,06; 0,64] & 92,3 (80,7–96,4) & 85,0 (82,7–86,7) & [0,23; 0,97; 0,03; 0,01] & +1,3 \\
R$^4$ — 20-shot† & 90,4 (86,9–95,4) & 68/62 & 94,7 (93,1–97,7) & 86,5 (84,7–88,0) & [0,15; 0,76; 0,06; 0,63] & 92,9 (90,0–96,2) & 85,3 (83,3–88,0) & [0,23; 0,97; 0,03; 0,01] & +0,9 \\
R$^4$ — 40-shot† & 91,6 (86,4–94,5) & 58/52 & 93,1 (90,0–95,5) & 86,7 (84,7–88,7) & [0,16; 0,75; 0,03; 0,64] & 91,3 (88,2–93,6) & 85,3 (84,7–86,0) & [0,21; 0,98; 0,02; 0,00] & +1,0 \\
R$^4$ — W ajustado ao rótulo real (teto)† & — & — & — & 87,3 (86,7–88,7) & [0,15; 0,71; 0,03; 0,68] & — & 86,4 (86,0–86,7) & [0,22; 0,98; 0,00; 0,00] & — \\
\bottomrule
\multicolumn{10}{p{0.95\linewidth}}{\footnotesize Meia-lua (150 pontos, sem split). O LLM vê sempre só x1, x2; R$^3$ acrescenta x1·x2 e R$^4$ acrescenta x1², x2² apenas na estimação da métrica. Colunas e convenções da bloco1\_table\_01: o teto (W ajustado ao rótulo real) só tem acurácia (vs rótulo real), não fidelidade. Passos do Problema D: zero-shot = Passo 2 (o LLM rotula sem exemplos); $k$-shot = Passo 3 (o LLM rotula com $k$ exemplos do rótulo real); a métrica de cada linha é estimada no Passo 4; teto = Passo 5 (W ajustado diretamente ao rótulo real). Em zero-shot (e com nomes de classe neutros) a correspondência entre nome e classe é arbitrária: a acurácia (vs rótulo real) do LLM pode sair abaixo de 50\%; o valor invariante à permutação é max($a$, 1 $-$ $a$). † = em pelo menos uma estimativa o Perceptron não encontrou $\hat{W}$ que separe os rótulos: melhor esforço ($\gamma = 0$) ou margem bruta mínima negativa (algum ponto do lado errado da fronteira). ``—'' = colapso em classe única (ou quase-colapso: classe minoritária com menos de 5 pontos); na linha colapsada também a acurácia do LLM e a partição ficam como ``—''.} \\
\end{tabular}
\end{table}
